% E-from-v1.tex: что версия 2 сохранила от версии 1 и что исправила; версия 3 сохраняет то же.
\section{От версии 1 к версии 2}
\label{app:from-v1}

Версия 1 остаётся историческим отчётом и не правится. Настоящее приложение записывает, что
версия 2 сохранила, что исправила и что сняла, чтобы читатель обеих мог их согласовать. Версия 3
сохраняет каждое решение ниже без изменений; то, что она изменила по сравнению с версией 2, есть
приложение~\ref{app:from-v2}.

\subsection*{Сохранено}

\begin{itemize}
  \item Геометрия сущностей с токеном в центре (рис.~\ref{fig:nucleus}), воспроизведённая из
    версии 1 без изменений и используемая как ядро каждой карты настоящего документа.
  \item Автомат жизненного цикла (рис.~\ref{fig:lifecycle}), воспроизведённый без изменений.
  \item Довод об уровне раскрытия: уровень есть сохраняемое свойство записи о проверке,
    закреплённое ограничением целостности, и именно это не даёт журналу проверок служить базой
    данных для слежки.
  \item Довод о том, что постквантовой подписи место с первого дня, через горизонт секретности
    идентификационных данных и неравенство Моски, а также довод об основании, согласно которому
    всякое высшее свойство идентификационной системы производно от её криптографических
    примитивов. Оба кратко изложены здесь и полностью стоят в версии 1; версия 3 сужает
    \emph{утверждение} о подписях, а не довод в пользу их выбора (раздел~\ref{sec:signatures}).
  \item Сравнение проектных свойств с развёрнутыми системами, которого настоящий документ не
    повторяет; первая страница репозитория несёт текущую таблицу и помечает Полярис как не
    развёрнутый и не выдающий удостоверений в национальном масштабе.
  \item Рамка модели угроз по шести классам угроз, которую проектная запись репозитория теперь
    несёт так, что каждая угроза сопоставлена средству защиты или явно принятому риску.
\end{itemize}

\subsection*{Исправлено}

\begin{itemize}
  \item Двенадцать таблиц стали тридцатью шестью в версии 2 и сорока пятью в версии 3. Три
    вспомогательные сущности версии 1 (привязки устройств, якоря реестра, список отзывов)
    остаются; остальной рост описан кольцо за кольцом в разделе~\ref{sec:walls}.
  \item Версия 1 утверждала, что таблицы аудита работают только на добавление по соглашению и
    благодаря инструментам и что автомат состояний будет закреплён триггером в промышленном
    развёртывании. Теперь и то и другое закреплено триггерами, и проверка роняет сборку, если хоть
    одно из них перестанет действовать.
  \item Версия 1 описывала якорь реестра как необязательный путь децентрализованных
    идентификаторов. Обязательство пакетов вне цепи существует и находится под наблюдением как
    журнал прозрачности; публикация во внешнюю цепь остаётся действием оператора с объявленными
    драйверами и без живой интеграции.
  \item Версия 1 моделировала семь контекстов проверки с биометрическими требованиями для каждого
    контекста и физический токен с биометрической привязкой на устройстве. Контексты остаются;
    физический токен теперь есть кодировка, эмулятор и устройство проверки и остаётся
    непроизведённым.
  \item Два умозрительных горизонта привязки из последнего приложения версии 1 (генетическая
    привязка и привязка к квантовому наблюдателю) не переносятся как проект. Проход репозитория по
    готовности отправил умозрительную схему в карантин за проверкой, не позволяющей ей вернуться, и
    настоящий документ считает расширением, которое нужно будущей среде, ключ на стороне
    держателя, а не новую биометрию.
\end{itemize}

\subsection*{Добавлено}

Всё, начиная с раздела~\ref{sec:signatures}, описывает механизмы, которых не существовало, когда
писалась версия 1: настоящее подписание и хранение ключей, двух свидетелей, отделяемый верификатор
и контракт соответствия, состояние и автономные полномочия, доказательство принадлежности и его
второго свидетеля, протокол кошелька и брокер, федерацию, прозрачность, ткань обмена,
эксплуатационные учения, Атлас и Афину, историю самоисправления, общественные границы и
обращённые в будущее разделы об агентах, доказательстве того, что ты человек, передаче полномочий
и внешних наблюдателях.
